\section{单圈紫外发散}
\label{sec:UVOL}
我们在图~\ref{fig:oneloop-g}中给出动量为 $p$ 的渐近胶子的准胶子算符相关的一圈费曼图，其中图 (e) 中的泡泡包含了主动胶子的所有一圈自能图。要得到完整的一圈贡献，还需要额外的费曼图：其中一些是图 (b)、(c)、(d)、(e) 和 (g) 的镜像，其余则可通过在这些费曼振幅中把外动量 $p$ 换成 $-p$ 得到。在下面的一圈计算中，我们以文献~\cite{Ji:2013dva}中的线性组合准胶子算符为例，但我们的结论对 36 个独立算符中的任何一个都成立。

%%%%%%%%%%%%%%%%%%%%%%%%%%%%%%%
\begin{figure}[htb]
\begin{center}
\includegraphics[width=0.75\textwidth]{images/oneloop-g.pdf}
\caption{\label{fig:oneloop-g}
动量为 $p$ 的渐近胶子的准胶子 PDF 在一圈阶的一些典型费曼图。}
\end{center}
\end{figure}
%%%%%%%%%%%%%%%%%%%%%%%%%%%%%%

我们取费曼规范，并为明确起见假设 $\xi_z$ 为正。图~\ref{fig:oneloop-g}中的图 (a) 给出
\begin{align}
M_{1{a}}&= \frac{g_s^2 \mu_r^{4-d} C_A}{i} \,e^{-i p_z \xi_z} \int_0^{\xi_z}dr_1\int_{r_1}^{\xi_z}dr_2 \frac{d^d l}{(2\pi)^d} \frac{e^{il_z(r_2-r_1)}}{l^2}
\nonumber\\
&\xlongequal{\text{UV}} \frac{\alpha_s C_A}{\pi} e^{-i p_z \xi_z} \left(-\frac{\pi \mu_r\xi_z}{3-d}+ \frac{1}{4-d} \right) \, ,
\label{eq:1a0}
\end{align}
其中 $\mu_r$ 是在维数正规化中补偿量纲的重正化标度。如预期的那样，该图同时对线性和对数紫外发散有贡献。

为了弄清式~$(\ref{eq:1a0})$ 中的紫外发散来自这一圈相空间何处，有益的做法是把圈动量 $l$ 的 $z$ 分量 $l_z$ 与 $l$ 的其余分量 $\bar{l}_\mu$ 区分开，即 $l_\mu = \bar{l}_\mu - l_z n_{\mu}$ 且 $l^2=\bar{l}^2-l_z^2$~\cite{Ishikawa:2017faj}。
如果在式~$(\ref{eq:1a0})$ 中把 $\bar{l}^2$ 约束在有限区域内，那么对 $l_z$、$r_1$ 和 $r_2$ 的积分不会产生任何紫外发散。此外，如果不包含 $|r_2-r_1|$ 非常小的区域，也不存在紫外发散；这可以通过引入如下分解来演示，
\begin{align}\label{eq:dec}
\frac{1}{(l+q)^2}
= \frac{1}{(\bar{l}+\bar{q})^2}+ \frac{(l_z+q_z)^2}{(l+q)^2(\bar{l}+\bar{q})^2}\, ,
\end{align}
其中 $q$ 可以是任何未积分的动量。把这个分解应用到式~\eqref{eq:1a0} 的 $\frac{1}{l^2} $ 上：第一项不含 $l_z$，因而对 $l_z$ 的积分给出 $\delta(r_2-r_1)$；第二项在对 $\bar{l}$ 积分后 UV 有限。也就是说，式~\eqref{eq:1a0} 中的紫外发散只能来自 $\bar{l}$ 处于紫外区而 $|r_2-r_1|$ 非常小的相空间区域。因此我们得出结论：在维数正规化下，图~\ref{fig:oneloop-g}中图 (a) 的全部紫外发散都来自一个定域于时空之中的区域。

利用式~\eqref{eq:dec} 同时分解 $\frac{1}{l^2} $ 和 $\frac{1}{(p-l)^2}$，我们对图 (b) 和 (c) 得到许多项，并发现这些项要么分母中不含 $l_z$（从而给出 $\delta(r)$ 及其导数），要么在对 $\bar{l}$ 积分后 UV 有限。因此这两个图的紫外发散同样定域于时空之中，
\begin{align}
M_{1{b}}
&\xlongequal{\text{UV}} \frac{\alpha_s C_A}{\pi} e^{-i p_z \xi_z} \left(\frac{-i}{p_z \xi_z} \frac{\pi \mu_r \xi_z}{3-d} \right) \, , \label{eq:1b0}\\
M_{1c}
&\xlongequal{\text{UV}} \frac{\alpha_s C_A}{\pi} e^{-i p_z \xi_z} \left(\frac{i}{p_z \xi_z} \frac{\pi \mu_r \xi_z}{3-d}+\frac{3}{4} \frac{1}{4-d} \right) \, ,
\label{eq:1c0}
\end{align}
其中图 (b) 只有线性紫外发散，而图 (c) 同时具有线性和对数紫外发散。

图~\ref{fig:oneloop-g}中的图 (d) 和 (e) 只有对数紫外发散，它们来自 $l^\mu$ 所有分量都趋于无穷的区域，因此也定域于时空之中。
图~\ref{fig:oneloop-g}中所有其他图形都没有紫外发散，原因很简单：由于 $\xi_z$ 有限，圈无法在时空局域化，这与准夸克算符的论证相同~\cite{Ishikawa:2017faj}。
综上，我们得出结论：准胶子算符在一圈的紫外发散只能出现在\textit{定域}于时空坐标中的区域。作为对比，定义 PDF 的算符的紫外发散来自沿“$-$”光锥方向\textit{非定域}的区域~\cite{Ishikawa:2017faj}。

来自图 (a) 的式~$(\ref{eq:1a0})$ 中的线性紫外发散是无害的：它可以像准夸克算符的情形一样，容易地被指数化为全阶的整体相位因子，然后被乘法重正化~\cite{Ishikawa:2017faj,Polyakov:1980ca,Dotsenko:1979wb}。然而，分别来自图 (b) 和 (c) 的式~$(\ref{eq:1b0})$ 与 $(\ref{eq:1c0})$ 中线性紫外发散的存在，可能挑战乘法可重正性。幸运的是我们发现，在维数正规化下这两个图的线性紫外发散彼此相消。相反，如果使用破坏规范对称性的截断正规化，图 (b) 和 (c) 的线性发散就不会相消~\cite{Wang:2017qyg,Wang:2017eel}。这意味着规范不变性在消除可能挑战乘法可重正性的线性紫外发散方面起着重要作用。下面我们将利用规范不变性证明：除规范链自能之外的所有线性发散在把全部贡献求和后都会相消，从而准胶子算符可以被乘法重正化。
